\chapter{CLI Command Reference}
\label{ch:appendix_cli}

This appendix provides a complete reference for all TriBench CLI command groups, their sub-commands, and arguments. The framework is operated entirely through this interface.

\begin{table}[H]
\centering
\small
\begin{tabular}{p{2.5cm}p{11.5cm}}
\hline
\textbf{Command} & \textbf{Description} \\
\hline
\multicolumn{2}{l}{\textbf{bundle} -- Bundle management commands} \\
\hline
\texttt{archive}   & Archive the bundle's results and logs \\
\texttt{clear}     & Clear the active bundle state \\
\texttt{create}    & Scaffold a new bundle directory \\
\texttt{info}      & Show bundle metadata and resolved paths \\
\texttt{list}      & List all experiment YAML files in the bundle \\
\texttt{set}       & Set a bundle as active persistently \\
\texttt{show}      & Show the currently active bundle \\
\texttt{validate}  & Check that the bundle has all required elements \\
\hline
\multicolumn{2}{l}{\textbf{sys} -- System lifecycle management commands} \\
\hline
\texttt{cluster}   & Manage the Kind Kubernetes cluster \\
\texttt{logs}      & Show system component logs \\
\texttt{port-forward} & Manage Kubernetes port forwarding tunnels \\
\texttt{setup}     & Set up the lakehouse system components \\
\texttt{start}     & Start the system and poll for readiness \\
\texttt{status}    & Check current system status \\
\texttt{stop}      & Stop the system gracefully \\
\texttt{teardown}  & Tear down the system (destructive operation) \\
\hline
\multicolumn{2}{l}{\textbf{exp} -- Experiment execution commands} \\
\hline
\texttt{cancel}    & Cancel a currently running experiment \\
\texttt{config}    & Show the configuration for a specific experiment \\
\texttt{list}      & List available experiments \\
\texttt{run}       & Execute an experiment end-to-end \\
\texttt{status}    & Check the status of a running experiment \\
\hline
\end{tabular}
\end{table}

\begin{table}[H]
\centering
\small
\begin{tabular}{p{2.5cm}p{11.5cm}}
\hline
\textbf{Command} & \textbf{Description} \\
\hline
\multicolumn{2}{l}{\textbf{res} -- Result viewing and basic analysis commands} \\
\hline
\texttt{analyze}   & Analyze experiment results (statistics, performance, scalability, etc.) \\
\texttt{archive}   & Archive old experiment results \\
\texttt{delete}    & Delete an experiment result by ID \\
\texttt{export}    & Export experiment results to CSV, JSON, or Parquet \\
\texttt{list}      & List all recorded experiment results \\
\texttt{monitoring} & Show hardware monitoring telemetry for an experiment \\
\texttt{queries}   & Show detailed per-query execution data \\
\texttt{reset-db}  & Reset the SQLite database by deleting all runs \\
\texttt{show}      & Show high-level summary of an experiment result \\
\texttt{suite-summary} & Show an aggregated summary across a suite of runs \\
\texttt{summary}   & Generate comparison summary for experiments \\
\hline
\multicolumn{2}{l}{\textbf{data} -- Dataset management commands} \\
\hline
\texttt{generate}  & Generate a benchmark dataset \\
\texttt{info}      & Show detailed dataset information \\
\texttt{list}      & List available datasets \\
\texttt{load}      & Load a dataset into the deployed system \\
\texttt{validate}  & Validate dataset integrity \\
\texttt{validate-iceberg} & Validate Iceberg tables in the catalog \\
\hline
\multicolumn{2}{l}{\textbf{config} -- Configuration management commands} \\
\hline
\texttt{defaults}  & Show framework defaults \\
\texttt{profile}   & Manage active configuration profile \\
\texttt{show}      & Show the current merged configuration \\
\texttt{trace}     & Trace where a specific configuration value originated \\
\texttt{validate}  & Validate the complete configuration for errors \\
\hline
\multicolumn{2}{l}{\textbf{suite} -- Experiment suite execution commands} \\
\hline
\texttt{list}      & List available experiment suites \\
\texttt{run}       & Execute all experiments in a designated suite \\
\texttt{show}      & Show detailed information about a suite \\
\hline
\end{tabular}
\caption{Complete reference of TriBench CLI commands, split across two tables.}
\label{tab:full_cli_reference}
\end{table}
